\chapter{Exact source entry for the strong HI corner}
\label{component:p54b8f20ad2dd}
\noindent\textbf{Component }\path{p54b8f20ad2dd}\par
{\small\noindent Source: \path{R13/baseline/baseline/advances/hi/exact_source_entry.tex}\par}
{\small\noindent SHA256: \path{54b8f20ad2dd194d0e2ed7f5124fa1cc688d91a27ff328d503dc53b0eb45bd9b}\par}
{\small\noindent Inventory use: Exact HI entry from n>=0, q,yInf and actual gaps; all four event weights unbounded\par}
\medskip
\section{Exact source entry for the strong HI corner}
Use the original positive rational C2 source with $d=1$, actual
$u_c$ unbounded and $v_c$ infinitesimal. Put
\[
 q=u_c^{-2},\quad y=v_c^2,\quad H=u_c^2h\in(0,1),
 \quad R=r/q,\quad B=b/q,\quad L=yt=B-H,\quad T=B-yR.
\]
The already proved source identities (WE1)--(WE3) show that $R$ is
bounded using only $n\geq0$, $q$ infinitesimal and $y$ bounded.
Here is their short contradiction argument, to make this use
explicit. The positive source gaps imply
\[
 0<T<H+y,\quad g=H+(1+y)r-1+qT>0,
\]
and $n=-qP/(2rQ)$ with $Q>0$ and
\[
 P\ge r^2(1+y)(H+yr)-Tqr(y+2)-q^2(T^2+T).
\]
If $q/r$ were infinitesimal, then $T/(H+yr)$ and
$(T^2+T)/(H+yr)$ would be bounded. For unit $r$ this follows
from $T<H+y\le(H+yr)/r$; for infinitesimal $r$, the inequality
$g>0$ forces $H$ to have residue one. Division by $r^2(H+yr)$
would give $P>0$, contrary to $n\geq0$. Thus $R$ is bounded.
(The exact positive remainder in this lower bound is printed in
(WE3), so this argument uses no asymptotic sign substitution.)

Now $r=qR$ is infinitesimal and $g>0$ gives $H_0=1$.
Define the positive opening
\[
 E=(1-H)/q.
\]
The same gap implies
$E<(1+y)R+T$, so $E$ is bounded. Since $j=1+R-t>0$,
$t$ is bounded; consequently
\[
 H=1-qE,\quad B=H+yt\sim1,\quad
 j=1+R-t,\quad K=B+R-E=g/q>0.
\]
All of $E,t,j,K$ are positive bounded, while
$t+j=1+R$ and $E+K=B+R$ are positive units.

Let $\Delta$ and $N$ denote the unchanged original source
polynomials, and set $\mathcal N=N/q$. Their exact expressions are
\begin{align*}
\Delta={}&(1+q)\{B^2+2BR+(1+y)B\}
 +(1-yq)R^2+2(1+y)R+1+y,\\
\mathcal N={}&(1+q)(B+R)(H+y)-(1+y)(1-qR)yt.
\end{align*}
Thus $\Delta$ is a positive bounded unit and
\[
 \bar\Delta=(1+\bar R)(3+\bar R),\qquad
 \overline{\mathcal N}=1+\bar R>0.
\]
Set $S=1+R+yt$, also a positive unit. Direct substitution into
the original coefficients gives the following exact equalities:
\[
\begin{aligned}
 mqK&=\frac{(1+y)(1-qR)(1+R)S}{2\Delta},&
 Cyt&=\frac{(1+R)\mathcal N}{2\Delta},\\
 \frac{A}{qy}&=\frac{(1+R)(B+R)S}{2H\Delta},&
 Dyj&=\frac{(1+q)(B+R)(B-yR)S}{2\Delta},\\
 \alpha qE&=\frac{(B+R)\mathcal N}{2\Delta},&
 n&=\frac{(B+R)(1+R)S-R\Delta}{2R\Delta}.
\end{aligned}                                                   \tag{HI-entry}
\]
Every one of the first five right-hand sides is a positive bounded
unit; in particular $B-yR$ has residue one. The last numerator has
residue $(1+\bar R)(1-\bar R)$, so $n\geq0$ implies
$0\le\bar R\le1$. This does not imply $R\le1$ as a source
inequality, and no such stronger claim is used.

Since $q,y$ are infinitesimal and $E,t,j,K$ bounded, the positive
numbers $qK,yt,yj,qE$ are infinitesimal. Dividing their unit
products in (HI-entry) proves
\[
 m,C,D,\alpha\quad\hbox{are all unbounded},\qquad
 A\asymp qy.
\]
This conclusion follows from the original source formulas and
$n\geq0$; no additional strong-weight assumption is needed for
an actual HI source.
